\chapter{Параметрические и непараметрические критерии}
\label{ch:hypothesis}

В главе~\ref{ch:basics} введена общая схема проверки гипотез: нулевая гипотеза $H_0$, альтернатива $H_1$, статистика $T$ с известным нулевым распределением и p-value как мера экстремальности наблюдения при $H_0$.
Там же показано, что выбор критерия начинается не с формулы, а с постановки вопроса: что именно считается отсутствием эффекта, какие наблюдения независимы, какая шкала у переменной и насколько данные по форме согласуются с предполагаемой моделью.
Настоящая глава развивает этот аппарат в двух направлениях, которые в прикладной работе дополняют друг друга, а не конкурируют.

В \S\ref{sec:parametric} рассматриваются параметрические критерии: предполагается, что выборка порождена распределением из заданного семейства (нормальное, Бернулли), а $H_0$ формулируется через неизвестные параметры.
Если предпосылки разумны, такие процедуры обычно обладают наибольшей мощностью: модель узкая, зато статистика $T$ и её нулевое распределение выводятся точно или асимптотически.
В \S\ref{sec:nonparametric} --- непараметрические методы, которые минимально ограничивают форму распределения и опираются на знаки, ранги или перестановки наблюдений.
Их применяют, когда нормальность сомнительна, хвосты тяжёлые, шкала ординальная или когда $H_0$ касается медианы и всей функции распределения, а не только среднего.

Оба блока опираются на одну идею из гл.~\ref{ch:basics}: при справедливости $H_0$ распределение $T$ известно (точно или асимптотически), что позволяет количественно сопоставить данные с моделью.
Практический маршрут часто выглядит так: сначала визуализация и описательная статистика, затем параметрический критерий при обоснованных предпосылках, и при нарушении предпосылок --- непараметрический аналог или перестановочная проверка.
Примеры с крысами, временем ремонта линий связи и диаметрами шайб в этой главе проиллюстрируют обе ветви этой схемы.

\section{Параметрические критерии}
\label{sec:parametric}

Параметрический критерий применим, когда генеральная совокупность описывается семейством $\{F(\cdot;\theta):\theta\in\Theta\}$ \citep{lehmann2005}; гипотеза касается значения $\theta$ или его подмножества.
Типичные задачи --- сравнение среднего с нормативом, контроль дисперсии, сравнение центров двух групп --- всё это сводится к ограничениям на $\theta$ при фиксированной форме $F$.
Наиболее разработан случай нормальных выборок; для бинарных исходов естественной моделью служит распределение Бернулли.
В обоих случаях p-value вычисляется по той же логике, что и в гл.~\ref{ch:basics}: находим $t_{\text{obs}}=T(x^n)$ и оцениваем, насколько редки такие или более экстремальные значения при $H_0$.

\begin{definition}
Выборка $X^n=(X_1,\ldots,X_n)$ называется нормальной, если $X_i$ независимы и $X_i\sim N(\mu,\sigma^2)$ с неизвестными $\mu$ и $\sigma^2$.
\end{definition}

Для нормальных данных выводятся точные критерии Стьюдента, Фишера и $\chi^2$; при больших $n$ и выполнении условий ЦПТ аналогичные процедуры используют как приближение и для иных распределений.
Общий шаблон повторяется: стандартизируем разность между оценкой параметра и значением под $H_0$, делим на оценку стандартной ошибки и сравниваем полученное $t_{\text{obs}}$ с табличным или вычисленным нулевым распределением.
Различие между конкретными критериями --- в том, какой параметр проверяется и что считается известным ($\sigma$ или только выборка).
Единый маршрут главы можно пройти на одном учебном наборе: сначала Q--Q график и тест Шапиро--Уилка, затем $t$-критерий на среднее, при сомнении~--- Манна--Уитни, при необходимости~--- перестановки.
Если все четыре шага дают согласованный качественный вывод, p-value $t$-критерия интерпретируется увереннее; расхождение (как у времени ремонта линий)~--- сигнал проверить предпосылки, а не подбирать критерий апостериорно.

Z-критерий (известна $\sigma$): это частный случай схемы из гл.~\ref{ch:basics}, когда дисперсия генеральной совокупности задана априори (калиброванный прибор, длительный исторический контроль). Статистика имеет точное стандартное нормальное распределение при $H_0$: \[ H_0\colon \mu=\mu_0,\qquad Z=\frac{\bar X-\mu_0}{\sigma/\sqrt n}\sim N(0,1)\ \text{при}\ H_0.
\] Двусторонний p-value: $2\bigl(1-\Phi(|z_{\text{obs}}|)\bigr)$; односторонние~--- $F_{N(0,1)}(z)$ или $1-F_{N(0,1)}(z)$.
На практике $\sigma$ известна редко; Z-критерий полезен как иллюстрация и в задачах контроля качества с жёстко заданной точностью линии.

Одновыборочный $t$-критерий ($\sigma$ неизвестна): здесь дисперсию оценивают по выборке ($S$), и знаменатель случаен.
Именно поэтому $T$ следует распределению Стьюдента с $n-1$ степенями свободы, а не нормальному: \[ H_0\colon \mu=\mu_0,\qquad T=\frac{\bar X-\mu_0}{S/\sqrt n}\sim St(n-1)\ \text{при}\ H_0.
\] Двусторонний p-value: $2\bigl(1-F_{St(n-1)}(|t_{\text{obs}}|)\bigr)$.
При малых $n$ хвосты $St(\nu)$ тяжелее, чем у $N(0,1)$ (рис.~\ref{fig:z-t-null}), поэтому нельзя подменять $t$-критерий Z-критерием ради упрощения.

Критерий $\chi^2$ на дисперсию: проверяет не центр, а разброс. При нормальности $(n-1)S^2/\sigma^2$ имеет $\chi^2$-распределение: \[ H_0\colon \sigma^2=\sigma_0^2,\qquad \chi^2=\frac{(n-1)S^2}{\sigma_0^2}\sim\chi^2_{n-1}\ \text{при}\ H_0.
\]

Двухвыборочный $t$-критерий (независимые выборки, равные дисперсии): сравнивает средние двух групп при предположении $\sigma_1^2=\sigma_2^2$.
Объединённая оценка дисперсию $S_p^2$ повышает эффективность, но делает критерий уязвимым при её нарушении: \[ H_0\colon \mu_1=\mu_2,\qquad T=\frac{\bar X_1-\bar X_2}{S_p\sqrt{1/n_1+1/n_2}},\quad S_p^2=\frac{(n_1-1)S_1^2+(n_2-1)S_2^2}{n_1+n_2-2}, \] нулевое распределение~--- $St(n_1+n_2-2)$.

Критерий Уэлча (неравные дисперсии): знаменатель $S_1^2/n_1+S_2^2/n_2$, число степеней свободы~--- по Satterthwaite.
Его предпочитают по умолчанию при сравнении групп с разным разбросом (см. примеры с крысами и рабочими часами в этой главе).

Парный $t$-критерий (связанные наблюдения $D_i=X_{1i}-X_{2i}$): сводит пару измерений к одной выборке разностей; типичен для схемы до и после (когнитивные тесты при СДВГ): \[ H_0\colon \mathbb ED=0,\qquad T=\frac{\bar D}{S_D/\sqrt n}\sim St(n-1).
\]

F-критерий (равенство дисперсий): предшествует pooled $t$-критерию: если $\sigma_1^2\neq\sigma_2^2$, объединённая дисперсию вводит систематическую ошибку в тест на средние.
\[ H_0\colon \sigma_1^2=\sigma_2^2,\qquad F=\frac{S_1^2}{S_2^2}\sim F(n_1-1,n_2-1)\ \text{при}\ H_0.
\] Чувствителен к отклонениям от нормальности.
\begin{figure}[htbp]
    \centering
    \begin{subfigure}[b]{\FigWidth}
        \centering
        \includegraphics[width=\linewidth]{\FigDir/chapter-2/norm.png}
        \caption{$Z$-критерий: $N(0,1)$}
    \end{subfigure}\\[0.6em]
    \begin{subfigure}[b]{\FigWidth}
        \centering
        \includegraphics[width=\linewidth]{\FigDir/chapter-2/stud.png}
        \caption{$t$-критерий: $St(\nu)$}
    \end{subfigure}
    \caption{Нулевые распределения $Z$- и $t$-критериев; при малых $\nu$ хвосты $t$ тяжелее.}
    \label{fig:z-t-null}
\end{figure}

$t$-критерии чувствительны к выбросам и устойчиво работают при умеренных отклонениях от нормальности, если $n$ достаточно велико (центральная предельная теорема для $\bar X$).
При малых $n$ и заметной асимметрии следует проверять нормальность (график QQ, критерий Шапиро--Уилка) или переходить к непараметрическим аналогам (\S\ref{sec:nonparametric}) --- как мы сделаем для времени ремонта линий связи, где $t$-критерий окажется на грани значимости, а ранговый метод даст более устойчивый вывод.
Критерий Уэлча заменяет объединённую дисперсию $S_p^2$ на $S_1^2/n_1+S_2^2/n_2$ и использует поправку Satterthwaite для числа степеней свободы.
Одностороннюю альтернативу выбирают только если знак эффекта известен до получения данных; иначе p-value можно подогнать под желаемый вывод (пример с упаковкой).

Следующие примеры иллюстрируют одну схему из гл.~\ref{ch:basics}: формулируем $H_0$, вычисляем $T$, сопоставляем с нулевым распределением и интерпретируем $p$ содержательно, не отрывая числа от графика.
Пример: контроль веса упаковки (Z-критерий). На линии расфасовки норма среднего веса~--- 4~г при $\sigma=1$~г (дисперсия стабильна по технологическому регламенту).
Инспекция $n=9$ упаковок дала $\bar x=4{,}6$~г --- заметное, но не очевидное превышение на глаз.
\[ z=\frac{4{,}6-4}{1/\sqrt 9}=1{,}8,\quad p_{\text{двуст.}}\approx 0{,}072,\quad p_{\text{верх.}}\approx 0{,}036.
\] Двусторонняя альтернатива не отвергает $H_0$ при $\alpha=0{,}05$; односторонняя альтернатива о превышении нормы~--- отвергает. Это иллюстрирует опасность выбора альтернативы постфактум: та же выборка и та же $z_{\text{obs}}$, но разные $p$ и разные решения.

Пример: дисперсия объёма инъекции ($\chi^2$). Стандарт требует $\sigma^2=9$~кв.\,мл; при $n=25$ получено $s^2=12$.
Здесь $H_0$ касается разброса, а не среднего --- типичная постановка контроля воспроизводимости дозирования.
\[ \chi^2_{\text{obs}}=\frac{24\cdot 12}{9}\approx 32,\quad p_{\text{двуст.}}\approx 0{,}25 \] --- отклонения от стандарта статистически незначимы, хотя точечная оценка превышает норму. Содержательный вывод: наблюдаемое превышение может быть случайной флуктуацией при текущем объёме выборки.

Пример: программа ведения беременности ($t$-критерий). Для женщин за чертой бедности базовое среднее веса новорождённых~--- 2800~г. После программы при $n=25$: $\bar x=3075$~г, $s=500$~г. Одновыборочный $t$-критерий --- прямое применение схемы сравнения среднего с нормативом при неизвестной $\sigma$.
\[ t=\frac{3075-2800}{500/\sqrt{25}}=2{,}75,\quad p_{\text{двуст.}}\approx 0{,}011.
\] Программа статистически значимо связана с увеличением веса; размер эффекта $\approx 275$~г следует оценивать содержательно.
Содержательный вопрос здесь шире p-value: клинически значим ли прирост 275~г для новорождённого в данной популяции? Статистика указывает на отличие от нулевого флуктуационного шума; врач и эпидемиолог решают, достаточен ли он для рекомендации программы.
Такой разделённый вывод~--- норма прикладного анализа, а не противоречие статистике.

Пример: измерения скорости света. В классическом эксперименте Michelson зафиксировано $n=100$ измерений скорости света методом вращающегося зеркала --- редкий случай большой выборки при лабораторной точности.
Выборочное среднее $\bar x\approx 2{,}9985\times 10^5$~м/с, стандартное отклонение $s\approx 79$~м/с. На рис.~\ref{fig:speed-ttest} показано распределение наблюдений и 95\%-й доверительный интервал для $\mu$ (см. гл.~\ref{ch:basics}).

\begin{figure}[htbp]
    \centering
    \includegraphics[width=\FigWidth]{\FigDir/chapter-2/qqplot.png}
    \caption{Распределение измерений скорости света и интервал для среднего.}
    \label{fig:speed-ttest}
\end{figure}

Проверка $H_0\colon \mu=3{,}0000\times 10^5$ против двусторонней альтернативы даёт $t\approx -18{,}7$ и $p\approx 3\times 10^{-34}$: систематическое отклонение среднего от теоретического значения статистически значимо.
Содержательно это не опровержение физики, а сигнал о погрешности прибора или модели измерения; график подтверждает близость данных к нормальному распределению, что оправдывает $t$-критерий.
Важно: отклонение $H_0\colon\mu=\mu_0$ не означает отказ от нормальности --- отвергается конкретное значение среднего, а не форма $F$.

Пример: продолжительность жизни крыс. Зафиксирована продолжительность жизни крыс при двух режимах питания: restricted ($n_1=106$) и ad libitum ($n_2=89$).
Средние составляют $\bar x_1\approx 969$ и $\bar x_2\approx 684$~дней --- крупный содержательный сдвиг; вопрос в том, насколько он устойчив к предпосылкам $t$-критерия.

Двухвыборочный критерий Уэлча даёт $p\approx 3\times 10^{-16}$: ограничение калорийности ассоциировано с более длительной жизнью.
Ящики с усами на рис.~\ref{fig:mannwhitney} иллюстрируют сдвиг распределений; тот же вывод подтверждается ранговым критерием Манна--Уитни ($p\approx 8\times 10^{-17}$, см. \S\ref{sec:nonparametric}).
Этот пример связывает обе части главы: параметрический тест даёт убедительный результат, а непараметрический --- независимую проверку без предположения о нормальности сроков жизни.

Пример: рабочие часы (GSS, критерий Уэлча). Сравнение числа рабочих часов в неделю у работающих неполный день в 1974 ($n_1=108$) и 2014 ($n_2=196$) годах: среднее увеличилось примерно на 2{,}57~часа, $p\approx 0{,}027$ (критерий Уэлча).
Распределения асимметричны и с тяжёлыми хвостами; ящик с усами (рис.~\ref{fig:hpw}) полезен для интерпретации: медианы и квартили показывают сдвиг иначе, чем одно $\bar x$.

\begin{figure}[htbp]
    \centering
    \includegraphics[width=\FigWidth]{\FigDir/chapter-2/hpw.png}
    \caption{Рабочие часы в неделю: сравнение двух лет.}
    \label{fig:hpw}
\end{figure}

Пример: когнитивные тесты при СДВГ. В исследовании СДВГ содержатся результаты теста внимания ($D0$~--- до лечения, $D60$~--- через 60~дней) для $n=24$ детей.
Среднее изменение $\bar d\approx 4{,}96$ балла, медиана~--- $5$.
Парная структура данных обязательна: каждое измерение после сопоставлено с до у того же ребёнка у того же ребёнка.

Парный $t$-критерий для $H_0\colon \mathbb ED=0$ даёт $p\approx 0{,}004$: метилфенидат статистически значимо улучшает показатель.
Парный критерий знаковых рангов Уилкоксона даёт $p\approx 0{,}004$ --- вывод устойчив к отказу от нормальности разностей; мы вернёмся к парным ранговым процедурам в \S\ref{sec:nonparametric}.

Пример: сравнение дисперсий (F-критерий). При сравнении прочности керамики из двух партий по 220 изделий критерий Фишера на $H_0\colon\sigma_1^2=\sigma_2^2$ может дать $p\approx 0{,}17$ при визуально близких дисперсиях.
F-критерий особенно чувствителен к отклонениям от нормальности, поэтому незначимый результат F-критерия не лицензирует pooled $t$ без графика.
Перед применением pooled $t$-критерия имеет смысл проверить равенство дисперсий; при отклонении $H_0$~--- критерий Уэлча.

\begin{figure}[htbp]
    \centering
    \includegraphics[width=\FigWidth]{\FigDir/chapter-2/fish.png}
    \caption{Нулевое распределение F-критерия.}
    \label{fig:f-test}
\end{figure}

\label{sec:gof}
Многие параметрические процедуры предполагают нормальность.
Проверка $H_0\colon X\sim N(\mu,\sigma^2)$ --- частный случай критерия согласия (goodness-of-fit): $H_0$ задаёт полное распределение (возможно, с оцениваемыми параметрами).
Логика та же, что в гл.~\ref{ch:basics}: при $H_0$ статистика согласия имеет известное распределение; большие значения означают плохое совпадение данных с моделью.
На практике Q--Q график часто информативнее формального $p$: он показывает где именно модель расходится с выборкой (центр, хвосты, асимметрия).

Критерий Харке--Бера: ориентирован на асимметрию и тяжесть хвостов через моменты третьего и четвёртого порядка: \[ JB=\frac{n}{6}\left(\hat\gamma_1^2+\frac{\hat\gamma_2^2}{4}\right) \sim\chi^2_2\ \text{при}\ H_0, \] где $\hat\gamma_1$, $\hat\gamma_2$~--- выборочные асимметрия и эксцесс. Чувствителен к выбросам: один экстремальный $x_i$ может раздуть $\hat\gamma_2$.

Критерий согласия Пирсона ($\chi^2$ по ячейкам): классический подход для сгруппированных данных; разбиение области значений на $K$ интервалов, $n_i$~--- частоты, $p_i$~--- теоретические вероятности при $H_0$: \[ \chi^2=\sum_{i=1}^K\frac{(n_i-np_i)^2}{np_i}\sim\chi^2_{K-1-r}, \] $r$~--- число оценённых параметров. Требует $np_i\gtrsim 5$ в большинстве ячеек; иначе асимптотика $\chi^2$ ненадёжна.

Колмогорова--Смирнова (одновыборочный; Лиллиефорс~--- при подстановке $\hat\mu$, $\hat\sigma$): сравнивает эмпирическую и теоретическую функции распределения по максимальному вертикальному зазору: \[ D=\sup_x|F_n(x)-F_0(x)|; \] табличное нулевое распределение. Слабее чувствителен к хвостам, чем JB или Q--Q-диагностика на рис.~\ref{fig:qq-norm}.

Шапиро--Уилка: статистика $W$ на основе порядковых статистик и табличных коэффициентов; рекомендуется для формальной проверки нормальности при умеренных $n$, когда нужен одно числовое p-value дополнительно к графику.
\begin{figure}[htbp]
 \centering
 \begin{subfigure}[b]{\FigWidth}
 \centering
 \includegraphics[width=\linewidth]{\FigDir/chapter-2/qqplot.png}
 \caption{Нормальная выборка}
 \end{subfigure}\\[0.6em]
 \begin{subfigure}[b]{\FigWidth}
 \centering
 \includegraphics[width=\linewidth]{\FigDir/chapter-2/qq_tails.png}
 \caption{Отклонения на хвостах}
 \end{subfigure}
 \caption{Q--Q графики: согласие с нормалью и типичные отклонения на хвостах.}
 \label{fig:qq-norm}
\end{figure}

\begin{figure}[htbp]
 \centering
 \includegraphics[width=\FigWidth]{\FigDir/chapter-2/qq_skewness.png}
 \caption{Q--Q график при асимметрии выборки.}
 \label{fig:qq-skew}
\end{figure}

Рекомендуемый порядок действий:
\begin{enumerate}
 \item если данные явно дискретны или бинарны~--- специализированный метод;
 \item построить Q--Q plot; при умеренных отклонениях можно применять
устойчивые процедуры ($t$-критерий);
 \item для критериев, чувствительных к форме (F на дисперсии), при сомнениях~---
Шапиро--Уилк;
 \item на очень малых $n$ формальные критерии малоинформативны;
 на очень больших~--- отвергают даже ничтожные отклонения от нормальности.
\end{enumerate}
Незначимый тест на нормальность не доказывает идеальную нормальность --- лишь отсутствие статистически различимых отклонений при данном $n$.
При сомнениях разумно дублировать вывод непараметрическим критерием (\S\ref{sec:nonparametric}), как в примере с временем ремонта.
Важно различать два типа расхождения с нормальностью: формальный тест может не отвергнуть $H_0$ при малом $n$, хотя график уже показывает асимметрию; при большом $n$~--- наоборот, отвергает при ничтожном отклонении. Поэтому критерии согласия~--- не пропуск к $t$-тесту, а элемент сочетания графика и формального теста: сначала глаз, затем статистика, затем робастная проверка при содержательных сомнениях.

Пусть $\ell(\theta)=\log L(X^n;\theta)$~--- логарифм функции правдоподобия, $\hat\theta$~--- ОМП при отсутствии ограничений, $\theta_0$~--- значение под $H_0\colon \theta=\theta_0$.
Идея LR повторяет логику гл.~\ref{ch:basics}: сравниваем, насколько сильно ограниченная модель ($H_0$) проигрывает полной в подгонке к данным.
\begin{definition}
Статистика отношения правдоподобия (LR) равна
\[
\Lambda(X^n)=-2\log\frac{L(X^n;\theta_0)}{L(X^n;\hat\theta)} =-2\bigl(\ell(\theta_0)-\ell(\hat\theta)\bigr).
\]
\end{definition}
При стандартных регулярных условиях и справедливости $H_0$ асимптотически $\Lambda\sim\chi^2_r$, где $r$~--- число независимых ограничений в $H_0$.
LR-критерий универсален: для нормальной выборки проверка $\mu=\mu_0$ при известной $\sigma$ эквивалентна $Z$-критерию; при неизвестной $\sigma$~--- $t$-критерию.
Пример LR: для $H_0\colon\mu=\mu_0$ в нормальной модели $\Lambda$ монотонно связана с $t^2$; большое $\Lambda$ означает, что ограниченная модель с $\mu=\mu_0$ существенно хуже подгоняет данные, чем полная. Эта логика переносится на регрессию (гл.~\ref{ch:regression}), где LR-тест сравнивает вложенные GLM.
Преимущество LR~--- единая схема для сложных вложенных моделей (регрессия, GLM; см. гл.~\ref{ch:regression}).
На уровне этой главы LR связывает разрозненные $t$- и $\chi^2$-формулы в один принцип: чем больше $\Lambda$, тем менее правдоподобна $H_0$.

Для $X_i\sim Ber(p)$ i.i.d. типична проверка доли успехов: \[ H_0\colon p=p_0,\qquad \hat p=\bar X,\qquad \Lambda=-2\log\frac{L(p_0)}{L(\hat p)}\ \dot\sim\ \chi^2_1\ \text{при}\ H_0.
\] Точный биномиальный критерий: при $S=\sum X_i$ под $H_0$ $S\sim Bin(n,p_0)$; двусторонний p-value~--- сумма вероятностей реализаций не более вероятных, чем наблюдаемая.
Биномиальный критерий предпочтителен при малых $n$ и экстремальных $p_0$; LR и асимптотика нормального приближения $\hat p$ согласуются при $np_0(1-p_0)\gtrsim 5$.
Бинарная модель напоминает эксперимент с мышами и зеркалом (\S\ref{sec:nonparametric}): там $H_0$ сформулирована через медиану доли времени, а здесь --- через параметр $p$ при явной модели Бернулли.

\begin{itemize}
 \item Применять двухвыборочный $t$-критерий с объединённой дисперсией
 при явно различных $S_1^2$ и $S_2^2$ без проверки или перехода к критерию Уэлча.
 \item Интерпретировать значимый одновыборочный $t$-критерий как
 доказательство ненормальности~--- отклоняется конкретное значение $\mu_0$, а не форма распределения.
 \item Заменять биномиальный критерий нормальным при $n<20$
 и $p$ близком к $0$ или $1$.
 \item Считать незначимый тест на нормальность доказательством
 идеальной нормальности данных.
\end{itemize}

\section{Непараметрические критерии}
\label{sec:nonparametric}

Непараметрические критерии не предполагают конкретного семейства распределений; $H_0$ формулируется через медианы, функции распределения или симметрию.
Цена универсальности~--- меньшая мощность при выполнении параметрических предпосылок: если данные действительно нормальны, $t$-критерий использует информацию о величине отклонений, а знаковый критерий --- только их направление.
Зато при асимметрии, выбросах и ординальных шкалах непараметрические процедуры часто дают единственный честный p-value.
Примеры с крысами, временем ремонта и шайбами показывают, когда параметрический и непараметрический ответ совпадают, а когда расходятся.

\begin{definition}
Критерий знаков использует число наблюдений по одну сторону от порога (или число положительных разностей) при $H_0$ о медиане; при непрерывности нулевое распределение биномиальное.
\end{definition}

Одновыборочный критерий знаков ($H_0\colon \med X = m_0$): статистика $S^+=\#\{i\colon X_i>m_0\}$; при $H_0$ и непрерывном $F_X$ $S^+\sim Bin(n,1/2)$.
Интуиция: при верной $H_0$ о пороге $m_0$ наблюдения по разные стороны порога должны встречаться примерно поровну.

Парный критерий знаков ($H_0\colon \med D=0$ для $D_i=X_i-Y_i$): $S^+=\#\{i\colon D_i>0\}\sim Bin(n,1/2)$ при симметрии $D$ относительно нуля.
Критерий знаков применим при цензурированных данных (известен только знак отклонения от порога), ординальных шкалах и тяжёлых хвостах, где ранговые методы нестабильны.
Используется лишь информация о знаке, поэтому мощность меньше, чем у ранговых критериев, но предпосылки слабее.

Пример: предпочтение зеркала у мышей. В эксперименте с зеркалом у мышей приведены доли времени ($n=16$), проведённые мышью в камере с зеркалом --- непрерывные доли на $(0,1)$, но без явной нормальной модели.
Средняя доля $\bar x\approx 0{,}475$; проверка $H_0\colon \med X=0{,}5$ критерием знаков (число наблюдений $>0{,}5$ равно $6$) даёт $p\approx 0{,}79$: данных против симметрии вокруг $0{,}5$ недостаточно.
При этом $\bar x<0{,}5$, что иллюстрирует различие среднего и медианы (см. гл.~\ref{ch:basics}): критерий знаков отвечает на вопрос о медиане, а не о среднем арифметическом.
Этот пример --- мост от бернуллиевской логики доли успехов к непараметрической формулировке без жёсткой модели $Ber(p)$.

Критерий Уилкоксона (одна выборка, $H_0\colon \med X=m_0$): ранги $R_i=\operatorname{rank}|X_i-m_0|$, статистика $W^+=\sum_{i\colon X_i>m_0} R_i$; нулевое распределение~--- табличное или точное перебором $2^n$ знаков.

Парный Уилкоксон для связанных выборок~--- то же для $D_i$.

Критерий Манна--Уитни (две независимые выборки, $H_0\colon F_1=F_2$): $U_1=\#\{(i,j)\colon X_{1i}>X_{2j}\}$; при $H_0$ распределение $U_1$ таблично; для больших $n_1,n_2$ применяют нормальное приближение.
Важно помнить: $H_0$ здесь шире равенства средних --- это полное совпадение распределений; значимый результат может отражать и разную форму.
\begin{figure}[htbp]
 \centering
 \includegraphics[width=\FigWidth]{\FigDir/chapter-2/signrank.png}
 \caption{Нулевое распределение критерия знаковых рангов (иллюстрация).}
 \label{fig:signrank}
\end{figure}

Ранговые критерии требуют непрерывности (или мало совпадений) и для парного Уилкоксона~--- примерной симметрии распределения разностей при $H_0$.
Отклонение $H_0$ может отражать асимметрию, а не только сдвиг медианы; диагностика формы обязательна.
На рис.~\ref{fig:mannwhitney} сопоставлены распределения продолжительности жизни крыс (см. параметрический пример с крысами); ранговый критерий не предполагает нормальности и подтверждает различие групп тем же качественным выводом, что и критерий Уэлча, но другим путём вычисления $p$.

\begin{figure}[htbp]
 \centering
 \includegraphics[width=\FigWidth]{\FigDir/chapter-2/mann.png}
 \caption{Сравнение групп питания крыс: гистограммы и ящики с усами.}
 \label{fig:mannwhitney}
\end{figure}

Пример: время ремонта линий связи. Зафиксировано время устранения неисправности (в часах) для операторов ILEC ($n_1=1664$) и CLEC ($n_2=23$) --- типичная прикладная задача с сильной асимметрией и множеством нулей (мгновенное устранение).
Средние: $\bar x_{\text{ILEC}}\approx 8{,}4$, $\bar x_{\text{CLEC}}\approx 16{,}5$.
Критерий Уэлча даёт $p\approx 0{,}06$ (на границе $\alpha=0{,}05$), тогда как Манна--Уитни~--- $p\approx 0{,}001$.

Распределение времени ремонта сильно асимметрично (много нулей и длинный правый хвост), что нарушает предпосылки $t$-критерия; ранговый критерий устойчивее к выбросам и указывает на систематически более длительное устранение неисправностей у CLEC.
Здесь непараметрический вывод не противоречит параметрическому --- он уточняет его: при сомнительной нормальности и асимметрии опираться на $t$ при $\alpha=0{,}05$ рискованно, тогда как ранги дают устойчивый сигнал.
Ящик с усами (рис.~\ref{fig:boxplot}) показывает сдвиг медианы и хвостов нагляднее, чем одна пара $\bar x$.

\begin{figure}[htbp]
 \centering
 \includegraphics[width=\FigWidth]{\FigDir/chapter-2/boxplot.png}
 \caption{Ящик с усами: медиана, квартили и выбросы.}
 \label{fig:boxplot}
\end{figure}

Для проверки $H_0\colon F_1=F_2$ (полное совпадение распределений, не только средних) используют статистики, основанные на ЭФР.
Это ближе к духу гл.~\ref{ch:basics}: $H_0$ задаёт полную модель (здесь --- идентичность двух $F$), а $T$ измеряет расхождение данных с ней.

Критерий Смирнова (двухвыборочный Колмогоров--Смирнов): \[ D=\sup_x|F_{n_1}(x)-F_{n_2}(x)|; \] нулевое распределение табличное при $H_0$.

Критерий Андерсона (модификация Крамера--фон Мизеса для двух выборок): статистика на основе рангов объединённой выборки; чувствительнее к различиям в хвостах, чем Смирнов.
Отличие от Манна--Уитни: $H_0$ шире (равенство всей $F$), а не только сдвиг; при разных формах, но равных медианах, критерии согласия и ранговые критерии могут вести себя по-разному.
Для крыс и времени ремонта Манна--Уитни и $t$-критерий отвечали на вопрос о сдвиге; KS/Андерсон дополнили бы картину, если бы подозревали различие форм при близких медианах.

\begin{definition}
Перестановочный критерий строит нулевое распределение статистики $T$, перебирая перестановки или разбиения, равновероятные при $H_0$; p-value --- доля перестановок с не менее экстремальным $T$, чем наблюдаемое.
\end{definition}

Это тот же каркас, что в гл.~\ref{ch:basics}, но вместо табличного $F_T$ распределение $T$ симулируют из данных при верной $H_0$:
\begin{itemize}
 \item Одна выборка, $H_0\colon\mathbb EX=m_0$, симметрия $F$:
 $T=\sum(X_i-m_0)$; нулевое распределение~--- $2^n$ знаков.
 \item Связанные выборки: $D_i=X_{1i}-X_{2i}$, перебор знаков $D_i$.
 \item Независимые выборки, $H_0\colon F_1=F_2$:
 объединить выборки и всеми $\binom{n_1+n_2}{n_1}$ способами разбить на группы; $T=\bar X_1-\bar X_2$ (или функция рангов).
\end{itemize}
p-value~--- доля перестановок с $|T|\geq|t_{\text{obs}}|$ (или односторонне).
Выбор между перестановками, бутстрепом и табличным $t$-критерием зависит от того, что именно является $H_0$: равенство $F_1=F_2$, равенство медиан или только равенство средних при нормальности.
Перестановки наиболее строги по предпосылкам; бутстреп гибче, но не центрирует нулевое распределение при нарушении $F_1=F_2$.
На практике начинают с обоснованной параметрической процедуры и добавляют непараметрическую проверку, как в примерах с крысами и ADHD.
При большом $n_1+n_2$~--- случайная подвыборка из $10^4$--$10^5$ перестановок; СКО оценки $p$ порядка $\sqrt{p(1-p)/|G'|}$.
Перестановочный критерий при $H_0\colon F_1=F_2$ проверяет полное равенство распределений (сдвиг как частный случай), тогда как бутстреп для разности средних (гл.~\ref{ch:basics}) не требует равенства $F$, но даёт приближённый p-value. Перестановки~--- без возвращения; бутстреп~--- с возвращением; нулевое распределение перестановочного $T$ центрировано в значении при $H_0$, бутстреп-распределение~--- в $\hat\theta$.
Если предпосылки перестановочного критерия выполнены, он точнее и обычно мощнее бутстреп-аналога; бутстреп гибче при нарушении $F_1=F_2$.
\begin{figure}[htbp]
 \centering
 \includegraphics[width=\FigWidth]{\FigDir/chapter-2/ranksum10.png}
 \caption{Нулевое распределение ранговой статистики (иллюстрация).}
 \label{fig:ranksum}
\end{figure}

Пример: диаметры шайб (перестановки и ранги). При проверке соответствия среднего диаметра шайб нормативу ($H_0\colon \mathbb EX=m_0$ при симметрии) критерий знаковых рангов может дать $p\approx 0{,}067$, перестановочный критерий с $T=\sum(X_i-m_0)$~--- $p\approx 0{,}10$: оба не отвергают $H_0$ при $\alpha=0{,}05$, но используют разную информацию (ранги против суммы отклонений).
Это напоминает пример с мышами: при близком к границе $p$ выбор статистики и предпосылок меняет решение сильнее, чем кажется из одной цифры.
Для шайб параметрический $t$-критерий был бы уместен при нормальности; перестановки и ранги дают проверку без жёсткой формы $F$.

\begin{itemize}
 \item Считать непараметрический критерий полностью свободным от предпосылок: симметрия, независимость и связь $H_0$ с медианой или $F$ остаются обязательными.
 \item Выбирать Манна--Уитни как проверку разности средних:
 $H_0$ формулируется как равенство функций распределения; при разных формах, но равных медианах критерий может давать значимый результат.
 \item Игнорировать связность пар в задаче ADHD: двухвыборочный
 критерий вместо парного завышает ошибку I рода.
 \item Путать двухвыборочный KS (равенство $F$) с критерием на сдвиг средних.
\end{itemize}
При написании отчёта по сравнению двух групп полезна явная таблица: цепочку предпосылки, параметрический тест, робастный тест и вывод. Для крыс обе колонки совпадают; для ремонта линий~--- нет; для шайб все три процедуры близки к границе $\alpha$. Такая таблица снимает соблазн цитировать только удобное p-value.

Сопоставление критериев на одной задаче. Рассмотрим сравнение времени сборки детали (с) на двух линиях ($n_1=n_2=35$).
Ящики с усами показывают сдвиг медиан и несколько выбросов на линии~B.
Q--Q графики по группам: линия~A близка к нормали, линия~B~--- с лёгкой асимметрией. Левена даёт $p=0{,}12$; F на дисперсии~--- $p=0{,}18$: объединённый $t$ допустим, но критерий Уэлча~--- безопаснее по умолчанию.

Численно: $\bar x_A=12{,}4$, $\bar x_B=14{,}1$, Уэлч $p=0{,}04$, Манна--Уитни $p=0{,}02$, разность медиан $\approx 1{,}5$~с. Перестановочный тест с $T=\bar x_A-\bar x_B$ даёт $p=0{,}05$. Картина согласованная: линия~B медленнее; ранги чуть увереннее параметрического теста на границе $\alpha=0{,}05$.

В отчёте фиксируем цепочку: (1)~независимость наблюдений (разные детали); (2)~диагностика формы; (3)~$H_0\colon\mu_A=\mu_B$; (4)~$T$ и $p$; (5)~размер эффекта 1{,}7~с и его стоимость в производстве.
Если бы Q--Q показал сильную ненормальность, в тексте доминировали бы Манна--Уитни и перестановки, а $t$-критерий~--- как чувствительность, не как единственный аргумент. Этот шаблон связывает параметрический и непараметрический блоки главы в одну историю.

\section{Асимптотические критерии и критерии согласия для распределений}
\label{sec:theory-ch2}

Асимптотические критерии Вальда и меток (score). Для регулярного параметрического семейства $F(x,\theta)$, $\ell(\theta)=\log L(X^n;\theta)$, $\hat\theta_{\mathrm{MLE}}=\argmax_\theta\ell(\theta)$, $S(\theta)=\partial\ell/\partial\theta$, $I(\theta)=-\partial^2\ell/\partial\theta^2$.

\begin{critcard}{Критерий Вальда}
выборка: &
$X^n=(X_1,\ldots,X_n)$, $X\sim F(x,\theta)$ \\
нулевая гипотеза: &
$H_0\colon\theta=\theta_0$ \\
альтернатива: &
$H_1\colon\theta<\neq>\theta_0$ \\
статистика: &
$Z_W(X^n)=\dfrac{\hat\theta_{\mathrm{MLE}}-\theta_0}{\sqrt{\mathbb{D}\hat\theta_{\mathrm{MLE}}}}$
(на практике $\mathbb{D}\hat\theta_{\mathrm{MLE}}\approx I^{-1}(\hat\theta_{\mathrm{MLE}})$) \\
нулевое распределение: &
$Z_W\sim N(0,1)$ при $H_0$ \\
\end{critcard}
Достигаемый уровень значимости (piecewise): \[ p(Z_W)=\begin{cases} 1-F_{N(0,1)}(Z_W), & H_1\colon\theta>\theta_0,\\[4pt] F_{N(0,1)}(Z_W), & H_1\colon\theta<\theta_0,\\[4pt] 2\bigl(1-F_{N(0,1)}(|Z_W|)\bigr), & H_1\colon\theta\neq\theta_0.
\end{cases}
\] Критерий Вальда использует информацию о правдоподобии только в $\hat\theta_{\mathrm{MLE}}$; на конечных $n$ часто хуже LR и score.

\begin{critcard}{Критерий меток (score, Лагранжа)}
выборка: &
$X^n=(X_1,\ldots,X_n)$, $X\sim F(x,\theta)$ \\
нулевая гипотеза: &
$H_0\colon\theta=\theta_0$ \\
альтернатива: &
$H_1\colon\theta<\neq>\theta_0$ \\
статистика: &
$Z_S(X^n)=\dfrac{S(\theta_0)}{\sqrt{I(\theta_0)}}$ \\
нулевое распределение: &
$Z_S\sim N(0,1)$ при $H_0$ \\
\end{critcard}
\[ p(Z_S)=\begin{cases} 1-F_{N(0,1)}(Z_S), & H_1\colon\theta>\theta_0,\\[4pt] F_{N(0,1)}(Z_S), & H_1\colon\theta<\theta_0,\\[4pt] 2\bigl(1-F_{N(0,1)}(|Z_S|)\bigr), & H_1\colon\theta\neq\theta_0.
\end{cases}
\] Для $X_i\sim\mathrm{Ber}(p)$: $Z_S=\dfrac{\hat p-p_0}{\sqrt{p_0(1-p_0)/n}}$, $Z_W=\dfrac{\hat p-p_0}{\sqrt{\hat p(1-\hat p)/n}}$ (асимптотически эквивалентны LR).

$Z$-критерии для двух долей. Независимые выборки $X_1^{n_1}$, $X_2^{n_2}$~--- независимые, $X_k\sim\mathrm{Ber}(p_k)$.
В таблице $2\times 2$: в первой выборке $a$ единиц из $n_1$, во второй $b$ из $n_2$; $\hat p_1=a/n_1$, $\hat p_2=b/n_2$, объединённая оценка $P=(\hat p_1 n_1+\hat p_2 n_2)/(n_1+n_2)$.

\begin{critcard}{$Z$-критерий: разность двух долей (независимые)}
выборки: &
$X_1^{n_1}$, $X_2^{n_2}$ независимы, $X_k\sim\mathrm{Ber}(p_k)$ \\
нулевая гипотеза: &
$H_0\colon p_1=p_2$ \\
альтернатива: &
$H_1\colon p_1<\neq>p_2$ \\
статистика: &
$Z=\dfrac{\hat p_1-\hat p_2}{\sqrt{P(1-P)\bigl(\frac1{n_1}+\frac1{n_2}\bigr)}}$ \\
нулевое распределение: &
$N(0,1)$ \\
\end{critcard}
p-value~--- по той же кусочно-заданной схеме, что для $Z_W$.
Пример: $n_1=100$, $t_1=3$ дефекта; $n_2=200$, $t_2=5$~--- $H_0$ о равенстве долей дефектных не отвергается ($p\approx 0{,}8$).

Связанные выборки. Одни и те же $n$ объектов, два бинарных признака; таблица $2\times 2$: $e,h$~--- совпадения, $f,g$~--- несовпадения ($f$~--- несовпадения знаков признаков).

\begin{critcard}{$Z$-критерий: разность двух долей (связанные)}
выборки: &
$X_1^n$, $X_2^n$ связаны, $X_k\sim\mathrm{Ber}(p_k)$ \\
нулевая гипотеза: &
$H_0\colon p_1=p_2$ \\
альтернатива: &
$H_1\colon p_1<\neq>p_2$ \\
статистика: &
$Z=\dfrac{f-g}{\sqrt{f+g-(f-g)^2/n}}
=\dfrac{\hat p_1-\hat p_2}{\sqrt{\dfrac{f+g}{n^2}-\dfrac{(f-g)^2}{n^3}}}$ \\
нулевое распределение: &
$N(0,1)$ \\
\end{critcard}
Используют только дискордантные пары $(f,g)$, совпадения $e,h$ не входят в $Z$.
Пример (одобрение премьер-министра, $n=1600$): $e=794$, $f=150$, $g=86$, $h=570$; $H_0$ о неизменности рейтинга~--- $p\approx 2{,}8\times 10^{-5}$ (связанная схема), против $p\approx 0{,}022$ при ошибочном анализе как двух независимых выборок.

Критерий Ансари--Брэдли. Сравнение дисперсий при равных медианах (или после центрирования).
Объединённый вариационный ряд $X_{(1)}\le\cdots\le X_{(N)}$, $N=n_1+n_2$; ранги $\widetilde{\rank}$ присваивают от краёв к центру: $1,2,3,\ldots,3,2,1$.
$R_1=\sum_{i=1}^{n_1}\widetilde{\rank}(X_{1i})$~--- сумма таких рангов в первой группе.

\begin{critcard}{Критерий Ансари--Брэдли}
выборки: &
$X_1^{n_1}$, $X_2^{n_2}$ независимы; $\med X_1=\med X_2$ \\
нулевая гипотеза: &
$H_0\colon \mathbb{D}X_1=\mathbb{D}X_2$ \\
альтернатива: &
$H_1\colon \mathbb{D}X_1<\neq>\mathbb{D}X_2$ \\
статистика: &
$R_1(X_1^{n_1},X_2^{n_2})=\sum_{i=1}^{n_1}\widetilde{\rank}(X_{1i})$ \\
нулевое распределение: &
табличное при $H_0$ \\
\end{critcard}
Пример: вес свинцовых слитков у двух поставщиков, средние в норме; $H_0$ о равенстве дисперсий отвергается ($p\approx 0{,}014$).
 Перед применением проверьте предпосылку о центрах; иначе значимый результат может отражать сдвиг, а не разброс.

Критерии согласия по ЭФР (одна выборка). Семейство статистик строится на различии $F_n(x)$ и $F_0(x)$ (при оцениваемых параметрах нулевое распределение корректируют, как у Лиллиефорса для нормали).

Крамера--фон Мизеса (интеграл Смирнова).\[ W^2=\int_{-\infty}^{\infty}\bigl(F_n(x)-F(x)\bigr)^2\,dF(x).
 \]

Андерсона--Дарлинга.\[ A^2=\int_{-\infty}^{\infty} \frac{\bigl(F_n(x)-F(x)\bigr)^2}{F(x)\bigl(1-F(x)\bigr)}\,dF(x).
\] Вес $[F(1-F)]^{-1}$ усиливает вклад хвостов: при подозрении на тяжёлые хвосты AD предпочтительнее Колмогорова--Смирнова, чувствительного прежде к центру распределения.

\begin{critcard}{Критерий Андерсона--Дарлинга (одна выборка)}
выборка: &
$X^n=(X_1,\ldots,X_n)$ \\
нулевая гипотеза: &
$H_0\colon X\sim F_0$ (параметры $F_0$ фиксированы или оценены по выборке) \\
альтернатива: &
$H_1\colon H_0$ неверна \\
статистика: &
$A^2$ (интеграл выше; в ПО~--- по $F_n$ и $F_0$) \\
нулевое распределение: &
табличное / с поправкой при подстановке $\hat\mu$, $\hat\sigma$ \\
\end{critcard}
На практике AD дополняют Q--Q график: формальный $p$ при большом $n$ может отвергнуть ничтожные отклонения, при малом $n$~--- не заметить асимметрию, видную на графике.

Связь с другими критериями. Колмогоров: $D=\sup_x|F_n(x)-F(x)|$; критерий Пирсона~--- по ячейкам гистограммы ($\chi^2$). AD и $W^2$~--- промежуточный класс по чувствительности к форме; двухвыборочный критерий Андерсона (модификация $W^2$ по рангам объединённой выборки) проверяет $H_0\colon F_1=F_2$ целиком, а не только сдвиг медианы (см. \S\ref{sec:nonparametric}).

\section{Практика и семинар}
\label{sec:practice-ch2}
 Следующий блок повторяет семинар 2--3 курса ПСАД: задачи на закрепление теории, типичный порядок вычислений и интерпретации.
Каждая задача проходит три шага: (1)~ручной расчёт по формулам лекции; (2)~проверка в \texttt{scipy.stats} / \texttt{statsmodels}; (3)~содержательная интерпретация вместе с графиком. Фиксируйте $H_0$, $H_1$, $T_{\mathrm{obs}}$, $p$ и размер эффекта в единицах задачи.

Пудра Kanji (Z-критерий). Линия расфасовки: норма $\mu_0=4$~г, $\sigma=1$~г известно по регламенту; $n=9$ упаковок, $\bar x=4{,}6$~г. $H_0\colon\mu=\mu_0$, $H_1\colon\mu\neq\mu_0$; $z=(\bar x-\mu_0)/(\sigma/\sqrt n)=1{,}8$, $p_{\mathrm{двуст.}}\approx 0{,}072$.
Вывод: при $\alpha=0{,}05$ перегруз не доказан; односторонняя альтернатива о превышении нормы дала бы $p\approx 0{,}036$~--- иллюстрация опасности выбора альтернативы апостериорный.
Повторите через \texttt{statsmodels.stats.weightstats.ztest} на сэмплированной выборке с тем же $\bar x$ и $s$.

Дефектная партия (биномиальный критерий) $p_0=0{,}05$, $n=20$, $T=2$ дефекта.
$H_0\colon p\le p_0$, $H_1\colon p>p_0$; $p(T)=1-F_{\mathrm{Bin}(20,p_0)}(T-1)\approx 0{,}26$.
Вывод: превышения нормы по выборке не видно.
Сверка: \texttt{scipy.stats.binomtest} и \texttt{statsmodels.stats.proportion.binom\_test}.

Джеймс Бонд (мощность и ошибки) $12$ из $16$~--- взболтанный мартини.
$H_0\colon p=0{,}5$ против двусторонней/односторонней альтернатив.
Сравните \texttt{binom\_test}, \texttt{proportions\_ztest}, \texttt{proportions\_chisquare}.
Симуляция: для $n\in\{5,10,20,100,1000\}$ постройте кривые мощности при истинных $p$ и уровне $\alpha=0{,}05$; отдельно~--- эмпирическую ошибку I рода при $p_{\mathrm{ист}}=0{,}5$, $n=10$, перебирая $p_0$.
Вывод: малый $n$ даёт низкую мощность даже при $p=0{,}75$; Wald/z и точный биномиальный близки, но дискретность $T$ важна на малых $n$.

Опыт Майкельсона ($n=100$). Классические $100$ измерений скорости света: гистограмма, Q--Q plot, критерии Лиллиефорса (KS), Шапиро--Уилка, Харке--Бера.
Ожидание: распределение близко к нормальному; одновыборочный $t$-критерий на $\mu=\mu_0$ даёт крайне малый $p$ из-за большого $n$, но это проверка среднего, а не идеальности формы.
Интерпретируйте отклонение $\mu_0$ как систематическую погрешность метода, а не как отказ от физики.

Разрушители легенд (доли, Spock). Часть 1 (тыльная сторона руки): $11$ из $12$ угадали свою фотографию среди $10$ чужих; $H_0\colon p=0{,}1$, $H_1\colon p>0{,}1$.
Биномиальный $p$ порядка $10^{-7}$; score-$Z$ и интервалом Уилсона согласуются.
Часть 2 (ладонь): $7$ из $12$; $p$ заметно больше, но часто $>0{,}05$ на малых $n$.
Сравнение долей на тыльной и ладонной сторонах как двух независимых $Z$-тестов (без таблицы связности) даёт $p\approx 0{,}03$~--- осторожная формулировка: тыл узнают лучше, но связанная схема была бы корректнее.
Аналог в лекции: дело Спока (90 женщин из 300 при $H_0\colon p=1/2$)~--- $p\sim 10^{-12}$.

Крысы (Уэлча, нормальность). Две группы: restricted и ad libitum; гистограммы, Шапиро--Уилка, Q--Q.
Усечение жизни $<400$ дней (смерть не от диеты) улучшает нормальность ($p_{\mathrm{Шапиро--Уилка}}\approx 0{,}05$ и $0{,}12$).
Уэлча на полных и усечённых данных: $p\ll 0{,}001$, разность средних $\approx 285$--$300$ дней.
Вывод: ограничение калорий связано с большей продолжительностью; ранговый контроль (гл.~\ref{ch:hypothesis}, пример с крысами) должен дать тот же качественный вывод.

Метилфенидат (парный $t$) $n=24$, парные разности $D_i$; Q--Q и Шапиро--Уилка на $D_i$.
Парный $t$-критерий: $p\approx 0{,}004$~--- эффект препарата значим. Вывод: средний прирост баллов положителен; без связности (ошибочный двухвыборочный $t$) выводы по $p$ и интервалу искажаются.
Дублируйте знаковым/ранговым критерием при отказе от нормальности.

Семинар~3: мыши и зеркало. Доли времени в камере с зеркалом ($n=16$); гистограмма.
$H_0\colon \med X=0{,}5$ или $\mathbb{E}X=0{,}5$ в зависимости от постановки.
Биномиальный тест ($6$ значений $>0{,}5$): $p\approx 0{,}79$.
Знаки, перестановки с $T=\sum(X_i-0{,}5)$, BCa-интервал для доли/среднего.
Вывод: данных против гипотезы о половине времени нет; $\bar x<0{,}5$ при неотвергаемой медиане~--- снова различие среднего и медианы.

Время ремонта Verizon (Манна--Уитни, KS). ILEC и CLEC: ящик с усами, гистограмма, Q--Q; много нулей и правый хвост. Уэлча: $p\approx 0{,}06$; Манна--Уитни: $p\approx 0{,}001$; перестановки и двухвыборочный KS согласуются с рангами.
Вывод: у CLEC систематически дольше устранение неисправностей; на границе $\alpha$ доверять только $t$-критерию рискованно.

Дополнительные задачи sem3 (кратко). Анорексия: парные веса; при отклонении от нормальности~--- знаки, Уилкоксон, перестановки; BCa-интервал прироста веса.
Недвижимость Сиэттл: $n_1=n_2=50$; $t$, Манна--Уитни, перестановки; бутстреп-интервалы средних $[246,330]$ и $[241,403]$, разности $[-52,115]$~--- сдвиг цен неочевиден.
Алюминий в тополях: $13$ клонов, парные концентрации; знаки, Уилкоксон, перестановки ($p\approx 0{,}03$ при $T=t$); BCa для разности сезонов.

\section{Контрольные вопросы по теме}

\begin{enumerate}
    \item Сформулируйте $H_0$ и $H_1$ для одновыборочного $t$-критерия
          в примере со скоростью света. Почему отклонение $H_0$ не означает, что данные ненормальны?
    \item Чем Z-критерий отличается от $t$-критерия? Когда различие существенно?
    \item Запишите статистику $\chi^2$ для проверки $\sigma^2=\sigma_0^2$
          и укажите нулевое распределение.
    \item Чем критерий Уэлча отличается от классического двухвыборочного
          $t$-критерия?
    \item Сравните JB, KS (Лиллиефорс) и Шапиро--Уилка: чувствительность
          и ограничения.
    \item Запишите статистику LR для $H_0\colon p=p_0$ в модели Бернулли.
    \item Почему критерий знаков для эксперимента с зеркалом у мышей не отвергает
          $H_0\colon \med X=0{,}5$, хотя $\bar x<0{,}5$?
    \item Сравните выводы $t$-критерия Уэлча и Манна--Уитни для
          задержек Verizon.
    \item Опишите алгоритм перестановочного p-value для независимых выборок.
          Чем он отличается от бутстрепа?
    \item Чем двухвыборочный критерий Смирнова отличается от Манна--Уитни?
    \item Верно ли, что «p-value $=0{,}03$ означает, что $H_0$ верна с вероятностью 3\%»? Объясните.
\end{enumerate}
